% TOP_PROBLEM: {"id": 1100307, "title": "Questions in Geometric Group Theory — Q 3.7", "category_id": 17, "category": {"id": 17, "name": "group_theory", "display_name": "Group Theory", "description": "Problems about groups, group actions, representations, and related algebraic structures.", "slug": "group-theory", "order_index": 17, "created_at": "2026-07-31T15:26:25.671Z"}, "status": "open", "problem_number": "AMR-010-0307", "set_id": 13}
% FIRST_POSTED: 2026-10-01
% ENTRY_KIND: statement_only
% UnsolvedMath ID 1100307; source catalogue code AMR-010-0307
% !TeX program = lualatex
% Definition and sources only; this file is not an LLM solution attempt.
\documentclass[11pt,a4paper]{article}
\usepackage[margin=25mm]{geometry}
\usepackage{fontspec}
\setmainfont{lmroman10-regular.otf}[BoldFont=lmroman10-bold.otf,
 ItalicFont=lmroman10-italic.otf,BoldItalicFont=lmroman10-bolditalic.otf]
\usepackage{amsmath,amssymb,mathtools}
\usepackage{unicode-math}
\setmathfont{latinmodern-math.otf}
\AtBeginDocument{\let\setminus\smallsetminus}
\usepackage{xurl,hyperref}
\hypersetup{colorlinks=true,urlcolor=blue}
\setcounter{secnumdepth}{0}
\setlength{\parindent}{0pt}
\setlength{\parskip}{5pt}
\setlength{\emergencystretch}{3em}
\begin{document}
\section*{Questions in Geometric Group Theory — Q 3.7}\label{1100307}
{\small UnsolvedMath ID 1100307 · AMR-010-0307 · Group Theory · AMR Open Problem Lists}
\subsection{Definitions and mathematical statement}
Let $F_m$ be the free group on $x_1,\ldots,x_m$, with $m\geq1$, and let $r\in F_m$ be a word. Write $\langle\!\langle r\rangle\!\rangle$ for its normal closure, the smallest normal subgroup containing it. The associated one-relator group is
\[
 G_r=F_m/\langle\!\langle r\rangle\!\rangle
      =\langle x_1,\ldots,x_m\mid r=1\rangle.
\]
The redundant case $r=1$ may be included; it simply gives a free group. A nontrivial relator can be taken cyclically reduced by cancellation and conjugation.

A limit group means a finitely generated fully residually free group. Explicitly, for every finite subset $E\subseteq G_r$, there must be a homomorphism $\phi:G_r\to F_2$ which is injective on $E$. Requiring only that each single nonidentity element survive in some free quotient would instead define ordinary residual freeness.

Characterize the one-relator groups that are limit groups. The target is a necessary and sufficient condition on the presentation, or an equivalent structural classification, valid for arbitrary finite rank and arbitrary relator. The criterion must be invariant under changes of free generators or other presentations of the same abstract group. It is not enough to identify individual families with surface-like relators or to test one selected finite subset. Proper powers, words involving fewer than all displayed generators, and presentations admitting free-product decompositions remain within the quantified class.
\subsection{Short English statement}
Characterize precisely which one-relator groups are fully residually free.
\subsection{Sources}
\begin{itemize}
\item[S1] ULAM.AI and the UnsolvedMath Contributors, supplied problems.json, record 1100307 (AMR-010-0307), AMR Open Problem Lists. Catalogue identity and source transcription; CC BY 4.0.
\url{https://huggingface.co/datasets/ulamai/UnsolvedMath}.
\item[S2] Bestvina - Questions in Geometric Group Theory (pdf) (2004), Question 3.7 (PDF page 11). Original source indexed by AMR-010-0307.
\url{https://www.math.utah.edu/~bestvina/eprints/questions-updated.pdf}.
\end{itemize}
\end{document}
